\documentclass[10pt,a4paper]{article} \usepackage[utf8]{inputenc} \usepackage[T1]{fontenc} \usepackage[polish,english]{babel} \usepackage[left=2cm,right=2cm,top=2cm,bottom=2cm]{geometry} \usepackage{enumitem} \usepackage{titlesec} \usepackage{hyperref} \usepackage{parskip} \hypersetup{ colorlinks=true, linkcolor=black, urlcolor=blue, pdfauthor={Anna Sheludko}, pdftitle={Anna Sheludko -- CV} } \titleformat{\section} {\normalfont\large\bfseries} {}{0em}{\MakeUppercase} [\vspace{0.2em}\hrule height 0.6pt] \titlespacing*{\section}{0pt}{1em}{0.5em} \pagestyle{empty} \setlist[itemize]{ left=1.2em, nosep, topsep=3pt, itemsep=2pt } \begin{document} \begin{center} {\LARGE\bfseries Anna Sheludko}\\[0.35em] Data Engineer \textbar{} Python Developer\\[0.3em] annasheludko152@gmail.com\\ \href{https://anna1sheludko.github.io}{anna1sheludko.github.io} \textbar{} \href{https://github.com/anna1sheludko}{github.com/anna1sheludko}\\[0.3em] Ukrainka \textbar{} mieszkam w Polsce \textbar{} dostępna do współpracy freelance \end{center} \section{Podsumowanie} Jestem inżynierką danych specjalizującą się w Pythonie, SQL oraz projektowaniu procesów ETL. Tworzę rozwiązania automatyzujące przetwarzanie danych i przekształcające rozproszone źródła w uporządkowane, gotowe do analizy bazy danych. Pracuję z Pythonem, PostgreSQL, Dockerem oraz narzędziami wykorzystywanymi do budowy pipeline'ów danych. Ostatnio zrealizowałam projekt ETL przetwarzający 1,5 mln rekordów w 9 tabelach, wyposażony w walidację schematu, automatyczne testy oraz wdrożony w środowisku Docker. \section{Umiejętności techniczne} \begin{itemize} \item \textbf{Języki programowania:} Python, SQL, PL/SQL, Java, C++, R \item \textbf{Bazy danych:} PostgreSQL, Oracle, MySQL, MS SQL Server, NoSQL, bazy semantyczne, grafy wiedzy \item \textbf{Przetwarzanie danych:} Pandas, NumPy, Pandera, Excel, CSV, JSON \item \textbf{Infrastruktura i DevOps:} Docker, Git, GitHub Actions, Linux, YAML \item \textbf{Testowanie i jakość:} pytest, walidacja danych, testy jednostkowe, dobre praktyki ETL \item \textbf{Dodatkowo:} REST API, podstawy Apache Airflow, statystyka, uczenie maszynowe, NLP, sieci komputerowe, technologie chmurowe \end{itemize} \section{Projekty} \textbf{End-to-End ETL Pipeline} \hfill \href{https://github.com/anna1sheludko/end-to-end-pipeline}{GitHub} \vspace{0.2em} \textit{Python, Pandas, PostgreSQL, Docker, Pandera, pytest} \begin{itemize} \item Zaprojektowałam i wdrożyłam kompletny pipeline ETL dla danych e-commerce obejmujący 9 tabel i około 1,5 mln rekordów. \item Zaimplementowałam mechanizm automatycznych ponowień podczas ekstrakcji danych, walidację schematu z wykorzystaniem Pandera oraz wydajne ładowanie danych przez PostgreSQL COPY, osiągając około 8-krotne przyspieszenie względem standardowych instrukcji INSERT. \item Przygotowałam konfigurację w YAML, strukturalne logowanie z rotacją plików, testy jednostkowe oraz automatyczne uruchamianie testów w GitHub Actions. \item Całość działa w pełni odtwarzalnym środowisku Docker Compose. \end{itemize} \vspace{0.6em} \textbf{API do PostgreSQL Pipeline} \hfill \textit{w realizacji} \vspace{0.2em} \textit{Python, REST API, PostgreSQL, Docker} \begin{itemize} \item Tworzę pipeline ETL pobierający dane z publicznego API, walidujący je i zapisujący w PostgreSQL na potrzeby dalszej analizy. \item Projekt obejmuje automatyzację pobierania danych, monitorowanie jakości oraz przygotowanie środowiska Docker. \end{itemize} \section{Wykształcenie} \textbf{Inżynieria danych (studia I stopnia)} \hfill \textit{planowane ukończenie: 2026}\\ Politechnika Świętokrzyska w Kielcach \begin{itemize} \item Kluczowe obszary kształcenia: bazy danych, hurtownie danych, bazy semantyczne, grafy wiedzy, uczenie maszynowe, statystyka, Python, Java, C++, sieci komputerowe oraz zarządzanie projektami. \item Sześciomiesięczne praktyki zawodowe realizowane w ramach programu studiów. \end{itemize} \section{Języki} \textbf{Ukraiński} (ojczysty) \textbar{} \textbf{Polski} (B2) \textbar{} \textbf{Angielski} (B2) \section{Dostępność} Jestem otwarta na współpracę przy projektach freelance, stażach oraz stanowiskach juniorskich związanych z Data Engineering i Python Development. Preferuję kontakt mailowy i zazwyczaj odpowiadam w ciągu 24 godzin. \end{document}